\section{Introduction}
\label{sec:intro}

Enterprises record regulations, failure analyses and operational lessons in
unstructured documents. Answering a question such as \emph{why did edge
delamination appear in a photovoltaic inspection?} requires finding the report
that records the decision, reading its evidence and citing it accurately.
Retrieval-augmented generation (RAG) addresses document discovery by embedding
queries and passages in a shared vector space \cite{rag-lewis,dpr}; agentic
systems extend this process by letting a language model call retrieval tools
\cite{mcp,singh-agentic-rag}. Yet semantic similarity measures topical
relatedness, not whether a passage contains evidence for the requested claim
\cite{siren,gao-ragsurvey}.

The relevant capabilities are usually treated separately. Retrieval asks
whether a relevant passage enters a ranked list. Knowledge management governs
how documents are organized and maintained. Document AI recovers text, tables
and formulas from complex layouts \cite{mineru}. In deployed systems, these
operations form one dependency chain: parsing errors affect indexing, placement
limits what retrieval can see, and ranking alone cannot establish whether a
candidate supports an answer. No single stage owns the property the user
actually needs---that the sentence being cited is supported by the document
being cited---and so no stage is accountable for it.

This paper presents QDCVR (Query-Driven Content-Verified Retrieval), a
self-hosted knowledge-base management system that treats content as the
shared object across acquisition, organization and question answering. It
organizes the workflow around three design contracts:

\begin{enumerate}
\item \emph{Content informs placement.} An ingest agent uses the parsed body
to choose a library and derive descriptions or tags; server-side checks
validate storage state and consistency, not semantic routing accuracy.
\item \emph{Content informs admission.} When the retrieval workflow invokes
the server-side Laya tool, it fetches and scores candidate segments and returns
survivors with provenance; the score is an admission signal, not truth.
\item \emph{Absence is reported.} The answer workflow instructs the agent to
report the searched scope when a completed evidence review finds no support.
Tool failures and incomplete searches are execution failures, not evidence of
semantic absence.
\end{enumerate}

The deployed platform integrates document parsing, content-informed
organization, lifecycle management, operational experiences, multi-lane
retrieval and a Model Context Protocol (MCP) server with 95 registered tools
\cite{mcp}. It runs locally and exposes retrieval functions to several agent
harnesses. Untested functions are identified rather than inferred from their
implementation.

\noindent\textbf{Contributions.}
\begin{enumerate}
\item A same-pool diagnostic that separates verifier \emph{ordering} from
verifier \emph{filtering}, and the transferable design conclusion it
supports. On identical logged industrial candidate pools, ordering by the
content-verifier score lowers Hit@5 from 0.94 to 0.44, while retaining fusion
order among recorded survivors restores 0.94; all 191 recorded scores exceed
the deployed threshold, so the loss is attributable to the induced ordering,
not to filtering. An independent NLI checkpoint fails the same way on the same
pools, and the collapse is not reproduced on a second, low-confuser corpus.
A verifier score that survives an admission filter is therefore not thereby
usable as a ranker, and the two roles have to be measured separately
(Sections~\ref{sec:lanes} and~\ref{sec:e4}--\ref{sec:e6}).
\item A knowledge-base architecture and workflow linking quality-gated
acquisition, content-informed placement, three-source consistency checks and
an invoked server-side evidence-admission tool, presented as three contracts
read as testable hypotheses rather than as achieved guarantees. For each
contract the mechanism is separated from the check that tests it, and the
evaluation reports the evidence level actually reached: placement (C1) is
checked mechanically and measured once without filename cues; admission (C2)
passed every recorded candidate on the industrial corpus, so its filtering
role went unexercised there and its ordering role is the failure in
contribution~1; abstention (C3) is answer-layer behavior observed in trials,
not a server-side classifier
(Sections~\ref{sec:system}--\ref{sec:implementation} and~\ref{sec:e8}).
\item A deployment evaluation across three corpora reporting failures and
measurement limits alongside successful checks, including unretrievable
papers with incomplete attribution, a reproduced catalog-transport failure,
and small-sample ablations without a detected end-to-end quality gain
(Section~\ref{sec:eval}).
\end{enumerate}

A CIKM~2026 demonstration introduced the organization and retrieval layers
\cite{qdcvr-demo}. The journal extension adds acquisition and lifecycle
measurements, the operational-knowledge case study, and the verifier-placement
and transport-failure analyses. The architecture and retrieval interface build
on the demo; the contribution claimed here is the additional evidence and
failure analysis, not a wholly new system.

\subsection{Research questions and reading guide}
\label{sec:rq}

Five questions structure the design-to-evidence sequence, ordered by the
weight of the answer rather than by the pipeline: whether a retrieved document
should be reordered by a verifier score or retain its fusion order, and what
happens when the verifier is applied under confuser density (RQ1); how content
commitments are implemented and checked across retrieval workflows and agent
interfaces (RQ2); how parsing checks, content-informed routing and lifecycle
audits behave in deployed collections (RQ3); how the tested retrieval
configurations compare on document recall, evidence markers and cost (RQ4);
and what abstention, latency, cost and execution-reliability observations
arise in end-to-end trials (RQ5).

Section~\ref{sec:related} positions the work; Sections~\ref{sec:system}--
\ref{sec:organization} describe the architecture and management functions;
Section~\ref{sec:lanes} specifies retrieval and verification;
Section~\ref{sec:implementation} covers deployment and enforcement; and
Section~\ref{sec:eval} evaluates the system, followed by implications,
limitations and conclusions in
Sections~\ref{sec:discussion}--\ref{sec:conclusion}.
